\documentclass[output=paper]{langscibook}
\ChapterDOI{10.5281/zenodo.22796683}
\author{Patrizia Giampieri\orcid{}\affiliation{Università degli Studi di Perugia}}
\title[The drafting style of employment agreements and employment offer letters]{The drafting style of employment agreements and employment offer letters: A corpus-based approach}
\abstract{Plain English is language that can be understood by the layperson. Over the decades, scholars and linguists have advocated the drafting of plain English texts in the legal field (\citealt{Tiersma1999}; \citealt{TiersmaSolan2012}; \citealt{Williams2023}). Nonetheless, many legal documents are still drawn up in the most intricate form, which is referred to as legalese. Even contracts addressed to non-native speakers of English (such as employment agreements) tend to be drafted in an archaic style. In this respect, there seems to be a significant difference between employment agreements and employment offer letters. This chapter is aimed at bringing to the surface the nuances between these two types of contracts. In particular, it unveils whether legalese features are used in both types of documents and if their drafting style is formal or informal. To this purpose, two corpora are constructed and analyzed: one composed of employment agreements, and the other of employment offer letters. This analysis highlights recurrent phrases, formulaic expressions, and legal terminology in such documents.

The chapter findings highlight that there seems to be a more pervasive use of plain English language in employment offer letters, whereas employment contracts tend to be drafted in legalese. The chapter presents implications of corpus-based analyses in translation studies. More precisely, it shows how corpus-driven activities can be adopted in the translation process. Its rigorous methods and principles may, in fact, be replicated in other fields, thus fostering interdisciplinarity in translation research.
\keywords{legal language, corpus analysis, legalese, plain English, the language of contracts, corpus-based translation}
}
\IfFileExists{../localcommands.tex}{
  \addbibresource{../localbibliography.bib}
  \usepackage{tabularx,multicol}
\usepackage{url}
\urlstyle{same}

 
\usepackage{langsci-optional}
\usepackage{langsci-lgr}
\usepackage{langsci-gb4e}

\usepackage{paralist}
\usepackage{enumitem}
\usepackage{cprotect}
\usepackage{fancyvrb}

\let\clipbox\relax
\usepackage[export]{adjustbox}
\usepackage{array}
\usepackage[normalem]{ulem}
\usepackage{langsci-branding}
\usepackage{langsci-bidi}
\usepackage[linguistics]{forest}

  
\makeatletter
\let\thetitle\@title
\let\theauthor\@author
\makeatother

\newcommand{\togglepaper}[1][0]{
   \bibliography{../localbibliography}
   \papernote{\scriptsize\normalfont
     \theauthor.
     \titleTemp.
     To appear in:
     E. Di Tor \& Herr Rausgeberin (ed.).
     Booktitle in localcommands.tex.
     Berlin: Language Science Press. [preliminary page numbering]
   }
   \pagenumbering{roman}
   \setcounter{chapter}{#1}
   \addtocounter{chapter}{-1}
}

\newbool{bookcompile}
\booltrue{bookcompile}
\newcommand{\bookorchapter}[2]{\ifbool{bookcompile}{#1}{#2}}

\newcommand{\qual}[1]{`#1'}
\newcommand{\bv}[1]{(\ref{#1})}

\newcommand{\ip}{ImPost}
\newcommand{\ipadv}{ImPost-adverbial}
\newcommand{\pradv}{ProxPost-adverbial}
\newcommand{\chisq}{Chi-squared}


% for rotated columns
\newcolumntype{T}[2]{%
    >{\adjustbox{angle=#1,lap=\width-(#2)}\bgroup}%
    l%
    <{\egroup}%
}
\newcommand*\rot{\multicolumn{1}{T{45}{1em}}}%


\VerbatimFootnotes


% \newcommand{\keywords}[1]{\textbf{#1}}


%  \newcommand{\textarabic}[1]{\RL{\arabicfont #1}}

  %% hyphenation points for line breaks
%% Normally, automatic hyphenation in LaTeX is very good
%% If a word is mis-hyphenated, add it to this file
%%
%% add information to TeX file before \begin{document} with:
%% %% hyphenation points for line breaks
%% Normally, automatic hyphenation in LaTeX is very good
%% If a word is mis-hyphenated, add it to this file
%%
%% add information to TeX file before \begin{document} with:
%% %% hyphenation points for line breaks
%% Normally, automatic hyphenation in LaTeX is very good
%% If a word is mis-hyphenated, add it to this file
%%
%% add information to TeX file before \begin{document} with:
%% \include{localhyphenation}
\hyphenation{
    par-a-digm
}

\hyphenation{
    par-a-digm
}

\hyphenation{
    par-a-digm
}

  \togglepaper[1]%%chapternumber
}{}

\begin{document}
\maketitle
%\shorttitlerunninghead{}%%use this for an abridged title in the page headers



\section{Introduction and research questions}\label{sec:giampieri:1}

This chapter is aimed at exploring the language of employment contracts and employment offer letters by unveiling differences in legal language and style. Exploring the linguistic similarities or differences in these two legal documents is relevant since they are recurrent in national and international contexts.

This chapter also wishes to scrutinize if and how corpora can be useful to carry out intralingual translations \citep{Jakobson1959}, i.e., whether they help translators produce texts which are plainer and/or more accessible.

The following questions are posed: 1) To what extent do these documents feature elements of \textit{legalese} or plain English?; 2) Are employment contracts and employment offer letters addressed to different categories of employees?, and 3) Do they differ in terms of content and purposes?

This chapter also discusses implications in the field of translation studies. Corpora, in fact, are useful when making hypotheses about language patterns in a second language. For this reason, further questions that this chapter wishes to address are the following: 4) Can corpora help discover the nuances of a second language for translation purposes?, and 5) Can corpora help translators deliver intralingual translations?

The added value of the study carried out in this chapter lies in the fact that it shows a valid corpus-based approach to understand and explore the language of employment legally binding documents. This contribution carries out detailed corpus-based analyses at both word and phrase level. Therefore, it shows possible translation strategies and tools which can be replicated in other academic studies and research. The novelty and usefulness of this approach are the rigorous methods applied and the way corpus consultation is carried out to bring to the fore language nuances and (dis)similarities between the two types of documents. Corpora provide linguistic resources that help users observe and analyse language change over time or in different documents or genres, also within the same domain. In this way, they showcase the linguistic diversity which translators can tap into when expressing terms and concepts in a target language.

\section{Theoretical background and literature review}\label{sec:giampieri:2}

Legal language is considered pedantic and archaic (\citealt{Bhatia1997}; \citealt{Williams2004}, \citealt{Williams2011} and \citealt{Williams2023}). Scholars claim that the “language of the law” \citep{Mellinkoff2004} is intricate for reasons of inclusiveness and preciseness (\citealt{CoulthardJohnson2007}: 40). Others argue that lawyers and judges do not have any particular communication intent; they use a complicate jargon only to remark their “legal fraternity”. \citep[52]{Tiersma1999}

\subsection{Legalese}\label{sec:giampieri:2.1}

Many legal documents are characterized by \textit{legalese} features which make it particularly abstruse, especially to the layperson (\citealt{Tiersma1999}: 200; \citealt{TiersmaSolan2012}). \textit{Legalese} features have been analyzed and discussed by many scholars over the years (\citealt{Bhatia1997}; \citealt{Tiersma1999}; \citealt{Williams2004}, \citealt{Williams2011}, \citealt{Williams2015} and \citealt{Williams2023}; \citealt{Giampieri2021}). Polysemy and near-synonyms, for example, characterize legal documents. Polysemous words are terms which have several meanings (\citealt{AlcarazVaróHughes2002}: 35ff), whereas near-synonyms are different words with very similar meanings (\citealt{Stubbs1996}: 109; \citealt{Philips2003}: 154). An example of polysemy is “consideration” which is a synonym of “payment” in legal discourse \citep[45]{Scott2017}, whereas in non-legal contexts, it means “reflection”, or “pondering”. The word “term” is also polysemic, as it can mean “duration”, “condition”, or even “word”, in certain circumstances. Hence, “term”, “duration” and “period” are near-synonyms in many legal contexts. Complex prepositions are another example of how legal language can be difficult to understand. Scholars refer to many instances in this respect, such as “by virtue of”, “for the purpose of”, “other than”, or “pursuant to\footnote{See also the list on \url{http://www.plainenglish.co.uk/the-a-z-of-alternative-words.html} and \url{https://www.plainlanguage.gov/guidelines/words/use-simple-words-phrases/}.} ” (\citealt{Bhatia2010}: 26–29; \citealt{CoulthardJohnson2010}: 10; \citealt{Mellinkoff2004}). As can be guessed, complex prepositions tend to make texts pedantic. Deixis also characterizes legal discourse, as readers have to move their attention back and forward in the text in order to understand what words refer to. Examples in this case are adverbial constructions such as “heretofore”, “thereon”, “thereto”, but also “aforementioned”, “foregoing”, etc. (\citealt{Abate1998}: 14–16). Ambiguity, amongst others, is achieved by resorting to “shall” as a modal verb to either express future actions or obligations, instead of using “will” or “must”, respectively (\citealt{Tiersma1999}; \citealt{TiersmaSolan2012}; \citealt{Williams2015} and \citealt{Williams2023}). Passive forms are very frequent and are used to “obscure the actors” \citep[75]{Tiersma1999}. Nominalization is another feature that hallmarks legal documents and conceals actors, or phrase subjects. With nominalization, verb and noun combinations are preferred over straightforward verbs (as in “make an agreement”, instead of “agree”) (\citealt{Bhatia2010}: 26–29; \citealt{Tiersma1999}: 75; \citealt{Williams2011}; \citealt{CoulthardJohnson2010}: 10). Also, in legal texts sentences are very long for reasons of all-inclusiveness (\citealt{Williams2004} and \citealt{Williams2015}). This gives way to syntactic discontinuity and embeddings (\citealt{Tiersma1999}: 56–57; \citealt{Williams2015}), which make the understanding of a passage particularly arduous. Finally, legal language tends to be sexist, given the predominance of male pronouns and male possessive adjectives \citep[135]{Griffith1988}.

\subsection{Plain English}\label{sec:giampieri:2.2}

In view of the “opacity” \citep{Orts2015} of the language of the law, over the years the Plain English movement has tried to make legal language simpler and more accessible to the layperson (see \citealt{Tiersma1999}; \citealt{Bhatia2010}; \citealt{TiersmaSolan2012}; \citealt{Williams2023}). Its advocates have proposed many strategies to simplify legal jargon by eliminating archaic expressions, reducing sentence length, using non-sexist language, removing unnecessary words, etc. (see \citealt{Bhatia2010}). Some institutions and drafters have embraced the principles of plain English in legal writing. For example, the  \citet{EnglishStyleGuide2023} developed at EU level is an innovative reference document. In the same way, the UK legislation has made huge strides forward and British drafters have produced more comprehensible acts and statutes \citep{Williams2023}.

\subsection{Legal language and corpora}\label{sec:giampieri:2.3}

Corpora are used to carry out language investigations \citep{Resi2024} and to foster second language learning, especially in technical fields (\citealt{chapters/5-Cuchet}). Given its specificity and technicalities, legal language has been studied and analyzed by using and consulting corpora (\citealt{VigierMorenoSánchezRamos2017}; \citealt{Biel2018};  \citealt{VigierMoreno2019}; \citealt{Giampieri2024}). Corpora help find and notice word usages in context, collocations and lexical phrases, which abound in legal texts. Legal language is so formulaic and repetitive that legal corpora need not be large to study and analyse its features (\citealt{BhatiaEtAl2004}: 207; \citealt{Giampieri2025}).

With regard to legal corpora, scholars posit that the ones publicly accessible are generally composed of court decisions and/or legislation (\citealt{Biel2018}: 29;  \citealt{PrietoRamosGuzmán2021}: 255). This means that no corpora on corporate or private legal documents are available due to confidentiality or privacy reasons (\citealt{VigierMorenoSánchezRamos2017}: 261; \citealt{Biel2018}). In this regard, however, the work by \citet{Giampieri2024} is innovative as corpora of terms of service are now publicly available. Generally, legal corpora on corporate documents tend to be compiled for specific research projects and are kept private (\citealt{VogelEtAl2017}: 12). Hence, studying the legal language of private legal texts must be made on a case-by-case basis by creating DIY (do-it-yourself) corpora \citep{Zanettin2002}. \citet[244]{Zanettin2002} suggests building a DIY corpus with 10–50 documents (5,000–40,000 words).

Corpora can be composed either automatically or manually. Automatic (or semi-automatic) corpus compilation is possible by using software solutions such as BootCaT (\citealt{BaroniBernardini2004}) or specific platforms, such as the Sketch Engine \citep{KilgarriffEtAl2004} (see, for example, the case studies by \citealt{Resi2024} and \citealt{Giampieri2024} and \citeyear{Giampieri2025}). Manual corpus creation is more time-consuming, but it is focused and targeted, as each document composing the corpus is carefully selected according to a specific project.

Corpora are consulted by translators as they help notice specific  language patterns (\citealt{Zanettin2002};    \citealt{VigierMorenoSánchezRamos2017}; \citealt{Giampieri2024} and and \citeyear{Giampieri2025}). They may also be exploited to assist translators in simplifying texts in intralingual translation processes \citep{StoddenEtAl2023}.

\section{Methodology and data collection}\label{sec:giampieri:3}

In order to investigate the legal language and style of employment contracts and employment offer letters, as well as unveil possible implications in translation studies, two specific corpora are built: a corpus of employment agreements and a corpus of employment offer letters. The corpora are created manually in order to make sure that the documents retrieved were drawn up by native speakers of English. Each document was, in fact, read carefully before considering it for the purposes of this analysis.

As for the employment agreement corpus, the following Google search string was written to retrieve employment contracts: \textit{“employment agreement{\textbar}contract” site:.onecle.com “governed by the laws of the State of”}. On the basis of this search string, it was possible to retrieve either “employment agreement” or “employment contract” in the onecle.com domain. Onecle.com is a database of US and non-US contracts. The string “governed by the laws of the State of” helped narrow down the results to US agreements. This is, in fact, a formulaic expression written at the end of North-American agreements stating the applicable law. The parties to the contract were also taken into consideration. To this aim, each webpage was read carefully to make sure that at least the drafting party was based in the US and, hence, was using native-speaking language. At the end of the process, 21 files were retrieved (3,812 word types; 104,633 tokens).

The Google search string of employment offer letters was the following: \textit{“employment offer letter” site:.onecle.com}. In this way, employment offer letters were downloaded from the onecle.com domain. As for the previous search, letters were carefully checked to verify that they had been written by native speakers of English. In this regard, only letters of employment issued by US companies were downloaded and saved. This process helped retrieve 31 files (2,650 word types; 46,235 tokens). As can be seen, despite being composed of more files, the second corpus shows fewer word types and tokens. This is due to the fact that employment offer letters are shorter than standard employment agreements. This peculiarity was noticed while composing the second corpus.

The corpora were analyzed (one at a time) by using the AntConc \citep{Anthony2024} software solution.

AntConc is an offline concordancer which lists concordances, collocations, KWiC (key words in context), and is equipped with many other functions. For the purposes of this study, the main search queries are used, such as concordancing and collocate listing.

This chapter carries out detailed analyses by following a methodological and rigorous approach. Hence, its methods and principles can be replicated in over fields of study, thus fostering interdisciplinarity.

\subsection{Data collection and analysis}\label{sec:giampieri:3.1}

The documents composing the corpora were analyzed firstly by following a general approach, then by adopting a more analytical perspective. In the first case, the texts of each corpus were skimmed in order to notice evident differences in their layout and general content. Afterwards, a detailed corpus-based analysis was carried out with the aim of revealing relevant language patterns. More precisely, \textit{legalese} features were searched for in both corpora, such as the modal auxiliary “shall”, passive constructions, deictic elements (such as “hereon” or “hereof”), polysemous words (e.g., “term” and “consideration”), etc.

\subsection{Feature selection}\label{sec:giampieri:3.2}

The \textit{legalese} features that were searched for in both corpora were chosen on the basis of literature findings. The use of “shall”, for example, has been criticized by many scholars (\citealt{Tiersma1999}; \citealt{Bhatia2010}; \citealt{Williams2023}) as it is an ambiguous auxiliary which can refer to both future actions and obligations (\citealt{Bhatia2010}: 28; \citealt{Williams2023}: 28–32). Sample phrases are “we shall see” and “shall be construed in accordance with” (\citealt{Williams2023}: 32–33). Passive forms are also used extensively in many legal documents to avoid mentioning an agent and to ensure the formality of statements \citep[32]{Williams2023}. An example is the sentence “the employment relationship shall be terminated upon the expiration of the legal term of notice”. \citep[32]{Williams2023} Deictic elements such as “hereon”, “hereto”, “hereof”, “thereon” etc. (\citealt{Abate1998}: 15; \citealt{Bhatia2010}: 26–28) are used to refer back and forward to parts of a document; in this way, they hinder smooth reading (e.g., “the parties hereto”, \citealt{Bhatia2010}: 16). Polysemous words such as “consideration” and “term” have been pointed out by the literature and criticized \citep{Scott2017}, as their meaning in daily usage is very different from the one they assume in legal contexts. For example “consideration” means “payment” in contracts \citep[45]{Scott2017}. Finally, simplified versions of complex prepositions such as “pursuant to”, “other than”, “for the purposes of”, and “by virtue of” have been proposed by the literature over the years (\citealt{Bhatia2010};  \citealt{ClearEnglish2014}). For instance, \citet{Bhatia2010} suggests using “under” or “according to” instead of “pursuant to” \citep[28]{Bhatia2010}, whereas the Clear English guide (\citeyear{ClearEnglish2014}) replaces “for the purpose of” with “to” \citep[2]{ClearEnglish2014}.

\section{Analysis}\label{sec:giampieri:4}

This section reports and discusses the corpus analysis of both corpora with a view to exploring the legal language and style of the employment agreements and employment offer letters.

\subsection{General observations}\label{sec:giampieri:4.1}

By scrutinizing the two types of documents from a general perspective, it is possible to notice some relevant differences.

The employment agreements composing the corpus, for example, are generally organized into clauses and start with premises. They state the duration of the contract; the duties and obligations of the parties; the compensation of the employee; the benefits; the terminating events; non-competition and non-disclosure clauses; dispute resolutions, and the governing law and jurisdiction. Standard “boilerplate” clauses (\citealt{Cordero-Moss2011}) are also mentioned, such as severability, entire agreement, waivers, etc.

Conversely, employment offer letters are not divided into clauses and have no recitals or premises at the beginning. Being letters, they address the reader/recipient by a “Dear” followed by the receiver’s first name (hence, their tone is quite informal since the very beginning). Texts tend to be organized into bullet points, reporting the employee’s start date, compensation, duties and benefits. Some refer to arbitration as a dispute resolution, whereas others offer stock options. Many of them express the hope for the recipient to accept the offer (e.g., “I hope that you elect to accept this offer of employment”) and highlight the sender’s availability in case of questions or doubts (e.g., “please feel free to call me if you have any questions at all”).

Given the general overview above, it can be noticed that employment agreements use a more formal tone and far more \textit{legalese} features than employment offer letters. As mentioned, many employment agreements start with premises (or recitals), whereas employment offer letters do not. Also, they feature opening formulae such as “Now, therefore, in consideration of the mutual covenants hereinafter set forth and for other good and valuable consideration, the parties hereto agree as follows:”, which the employment letters do not have. On the contrary, as already pointed out, letters of employment start with “Dear” followed by the employee’s first name and then they continue with phrases such as “we are / I am pleased to offer you the position of”, that make the reader feel at ease. Hence, there seems to be less detachment between the two parties. At the same time, employment agreements end with a closing formula, such as “In witness whereof, the parties have executed this Agreement as of”, which employment letters do not have. These, in fact, end in a much simpler way, with closing sentences such as “If you have any questions regarding this employment offer, please call (…)”, followed by the statement “I have read this offer and I understand and accept its terms”, or simply “agreed and acknowledged”. Furthermore, by opening and reading any employment offer letter, the use of the second-person pronoun “you” comes to the fore, as well as second-person adjectives (“your”). This makes the tone less formal and more direct (see also \citealt{Williams2023}). On the contrary, employment agreements use third-person references, such as “Employee” or “Executive”, and tend to use male pronouns and male possessive adjectives (“he”, “him”, “his”).

The findings above can be considered as relevant when translators need to simplify language and/or make it more accessible to laypeople (see \citealt{StoddenEtAl2023}). As a matter of fact, such observations reveal how intralingual translation strategies can be followed.

In light of these initial findings, it would be tempting to speculate that employment agreements are addressed to high-ranked managers, whereas employment offer letters to lower positions. This assumption is confuted by the fact that positions of Executive or Directors, for example, are offered in both types of documents. For instance, “Executive shall serve as executive vice president” is a phrase found in some employment agreements, and “offer for the position of Executive vice president” is mentioned in several employment offer letters.

\subsection{Corpus analysis of legalese patterns}\label{sec:giampieri:4.2}

An in-depth corpus analysis of the documents composing the two corpora is now carried out. The investigation compares the language used in the employment agreement (EA) corpus and in the employment offer letter (EOL) corpus, respectively. \tabref{tab:giampieri:1}  shows some small-scale quantitative data related to a comparative analysis of some \textit{legalese} features that is featured in both corpora. As mentioned, the \textit{legalese} patterns are chosen on the basis of literature findings (see \sectref{sec:giampieri:3.2} on \textit{legalese} feature selection). The occurrences of \textit{legalese} patterns are reported in entire figures and in percentage values. The percentage values are calculated against the number of word types of each corpus (i.e., 3,812 in the EA corpus and 2,650 in the EOL corpus).

\begin{table}
\begin{tabularx}{\textwidth}{l@{} rr Yr}
\lsptoprule
                    & \multicolumn{2}{c}{EA}  &\multicolumn{2}{c}{EOL}  \\
\cmidrule(lr){2-3}\cmidrule(lr){4-5}
{Legalese features} & {Count} & {\%} & {Count} & {\% }\\
\midrule
Shall & 1,559 & 40.9\% & 206 & 7.8\%\\
Passive forms & 1,082 & 28.4\% & 552 & 19.7\%\\
Deixis (e.g., “hereon”, “hereof”, “thereon”...) & 946 & 24.8\% & 127 & 4.8\%\\
Polysemous words (“term”) & 250 & 6.6\% & 43 & 1.6\%\\
Polysemous words (“consideration”) & 54 & 1.4\% & 4 & 0.1\%\\
Complex prepositions (“pursuant to”) & 139 & 3.6\% & 33 & 1.2\%\\
Complex prepositions (“other than”) & 97 & 2.5\% & 23 & 0.8\%\\
Complex prepositions (“for the purpose(s) of”) & 22 & 0.57\% & 1 & 0.03\%\\
Complex prepositions (“by virtue of”) & 4 & 0.1\% & 0 & {}-\\
\lspbottomrule
\end{tabularx}


\caption{A small-scale comparative quantitative analysis of legalese features in the employment agreement (EA) corpus and in the employment offer letter (EOL) corpus.}
\label{tab:giampieri:1}
\end{table}

Comments to the figures of \tabref{tab:giampieri:1} are reported on the next pages. The comments help understand the language used in both corpora. They also supplement the quantitative data of \tabref{tab:giampieri:1} by providing qualitative insights (i.e., sample phrases), and by presenting and discussing language patterning from the two corpora.

As far as “shall” is concerned, it is evident that the EA corpus makes a far larger use of this controversial modal verb (40.9\%). On the contrary, in the EOL corpus it is used more sparingly (7.8\%). Sample paragraphs sourced from the EA corpus are the following: “Employee shall devote all working time, attention, knowledge, and skills solely to the Company” and “Employer shall have no right to terminate under this Section”. Examples from the EOL corpus are the following sentences: “you shall (a) devote your entire working time for or at the direction of the Company” and “you shall not have the ability to terminate this Agreement”. \tabref{tab:giampieri:2}  reports these and other sentences.

\begin{table}
\begin{tabularx}{\textwidth}{lQQ}
\lsptoprule
{No.} & {Sample phrase (EA)} & {Sample phrase (EOL)}\\
\midrule
1 & \textit{Employee} shall devote all working time, attention, knowledge, and skills \textit{solely} to the Company. & \textit{You} shall (a) devote your entire working time for or at the direction of the Company.\\
2 & \textit{Employer} shall have no right to terminate under this Section. & \textit{You} shall not have the ability to terminate this Agreement\\
3 & \textit{Employer} shall reimburse \textit{Executive} for all traveling, hotel, entertainment and other out-of-pocket expenses & \textit{The} \textit{Company} shall promptly reimburse \textit{you} for all actual and reasonable business expenses.\\
\lspbottomrule
\end{tabularx}
\caption{Sample phrases with “shall”.}
\label{tab:giampieri:2}
\end{table}

Although “shall” is used in both corpora, it is observable that both EAs and EOLs use “will” to a different extent. The EOL corpus mentions it more systematically than the EA corpus (789 hits, 29.8\%, against 415, 10.9\%, in the EA corpus). The modal “must” also appears, but not considerably (34 hits, 1.3\%, in the EOL against 33 hits, 0.8\%, in the EA corpus).

In addition, the EA phrases reported in \tabref{tab:giampieri:2}   feature the words “Employee”, “Executive” and “Employer”, whereas the EOL statements use “you” and “The Company”. This aspect was already discussed and noticed above. Such appellatives highlight the fact that letters of employment tend to me more direct. Moreover, the EA corpus contains the formal term “solely” (line one), whose plain equivalent is “exclusively”, or “only”.\footnote{See \url{https://www.collinsdictionary.com/it/dizionario/inglese-sinonimi/solely}.}  In this respect, “solely” is used 17 times in the EA corpus (0.45\%), whereas only once in the EOL corpus (0.04\%).

Passive forms are not particularly easy to source via corpus analysis, given that the corpora are untagged. However, targeted searches can be carried out by writing the search string “be *ed”. By doing so, it is possible to retrieve passive forms composed of the bare infinitive “be” followed by a regular past participle. This kind of passive form is particularly frequent in the EA corpus (804 counts, 21\%, against 347, 13\%, in the EOL corpus). A similar search string, i.e., “is *ed”, confirms instances of passive constructs in the EA corpus (219 hits, 5.7\%, against 101, 3.8\%, in the EOL corpus). The search for “are *ed” shows slightly more hits in the EOL corpus (74, 2.8\%, against 59, 1.5\%, in the EA corpus). \tabref{tab:giampieri:3} exhibits corpus-sourced passive constructions.

\begin{table}
\begin{tabularx}{\textwidth}{lQQ}
\lsptoprule
{No.} & {Sample phrase (EA)} & {Sample phrase (EOL)}\\
\midrule
1 & Executive shall (…) and {he} shall \textit{not be required to relocate} his residence or principal office to any place outside (...) County. & \textit{You} will not \textit{be required to move} to the San Francisco Bay Area.\\
2 & \textit{It is agreed} by the parties \textit{hereto} that... & \textit{It is intended} that...\\
3 & Any prior agreements or understandings with respect to Employee’s employment by Employer, \textit{are canceled} as of the Effective Date. & Any contrary representations which may have been made to you \textit{are superseded} by this offer.\\
\lspbottomrule
\end{tabularx}
\caption{Sample phrases with passive forms.}
\label{tab:giampieri:3}
\end{table}

\tabref{tab:giampieri:3}   highlights that both corpora make use of passive forms. However, EOL clauses are less formal as they tend to feature plainer words and expressions. For example, instead of “relocate” (line one), the clause in the EOL corpus contains “move to”. Also, the first clause of the EA corpus uses sexist language (“\textit{he} shall”), whereas the corresponding clause of the EOL corpus displays “you”.

The search string “is *ed” (second line of \tabref{tab:giampieri:3}) reveals interesting patterns. For instance, the EAP corpus contains formal expressions such as “it is agreed”. This construct does not appear in the EOL corpus. Deixis also comes to the fore; in particular, with the adverbial expression “hereto” (line two of \tabref{tab:giampieri:3}).

As already mentioned, only the past participle of regular verbs is investigated. Nonetheless, it is evident that this type of passive form is more recurrent in the EA corpus (1,082 counts, 28.4\%, against 522 occurrences in the EOL corpus, corresponding to 19.7\%).

As noticed, deixis appears in some EA clauses. If “*here*” is searched for in both corpora, 946 hits are retrieved from the EA corpus (24.8\%), whereas only 127 (4.8\%) are found in the EOL corpus. In this regard, it should be highlighted that the “*here*” search provides also unrelated elements, such as “therefore”, “where” and “elsewhere” in the EOL corpus. Since “therefore” is a linking word and “where/elsewhere” are place adverbs, they are ruled out. The deixis elements obtained from both corpora are the following: “therein”, “thereof”, “thereto”, “thereafter”, “hereto”, “hereof”, “hereby”, “hereunder”, etc. \tabref{tab:giampieri:4} illustrates some examples.

\begin{table}
\begin{tabularx}{\textwidth}{lQQ}
\lsptoprule
 {No.} &  {Sample phrase (EA)} &  {Sample phrase (EOL)}\\
 \midrule
1 & Executive may sit on the boards of directors of corporations and committees
\textit{thereof} without violating his obligations \textit{hereunder}. & With respect to the subject matter \textit{hereof}.\\
2 & \textit{WHEREAS}, the Company and Employee entered into an Employment Agreement dated as of July 22, 1996 (as amended to the date \textit{hereof}, the “Employment Agreement”). & The Company consummates a change of control transaction, \textit{whereby} fifty percent (50\%) or more of the voting stock of the Company changes ownership.\\
3 & \textit{Thereafter}, this Agreement terminates and the Executive shall receive those benefits. & \textit{Thereafter}, you will accrue vacation at the regular (…) vacation accrual rate.\\
4 & The parties do \textit{hereby} agree as follow: [\textit{sic.}]. & You \textit{hereby} agree that the Confidentiality Agreement will govern and be applicable to your prior employment.\\
\lspbottomrule
\end{tabularx}
\caption{Sample phrases with deixis.}
\label{tab:giampieri:4}
\end{table}

Expressions such as “thereof”, “hereof”, “hereunder” (line 1) and “whereas” (line 2) or “thereafter” (line 3) make the reading and understanding of legal texts less smooth \citep{Williams2023}. As regards “hereby” (line 4), it could be argued that it is slightly more comprehensible or, at least, more frequent and accepted \citep[74]{Williams2023}. This adverb is featured in both corpora, although with different percentage values (84 hits in the EA corpus, 2.2\%, and 8 hits in the EOL corpus, 0.3\%).

The polysemous word “term” appears 250 times (6.6\%) in the EA corpus, and only 43 times (1.6\%) in the EOL corpus. With regard to the word “consideration”, the EA corpus produces more hits (1.4\% against 0.1\% of the EOL corpus). Sample paragraphs or phrases with such polysemous terms are reported in \tabref{tab:giampieri:5}.

\begin{table}
\begin{tabularx}{\textwidth}{lQQ}
\lsptoprule
 {No.} &  {Sample phrase (EA)} &  {Sample phrase (EOL)}\\
 \midrule
1 & The Original \textit{Term} and any 12 month period commencing immediately after the end of any \textit{Term.} & During the \textit{term} of your employment.\\
2 & The \textit{term} “Affiliate” shall include, but shall not be limited to... & The Severance Plan defines the \textit{term} “Cause” as a participant’s: Material neglect...\\
3 & (A) Upon all of the \textit{terms} and conditions described in this Agreement.\newline
(B) No waiver by any party of any \textit{term} or condition of this Appendix A & Under the \textit{terms} and conditions described above.\\
4 & As a condition of my employment with (…) its subsidiaries, affiliates, successors or assigns (together the “Company”), and \textit{in consideration of} my employment with the Company. & Like all Company employees, you were required, as a condition to and \textit{in consideration} \textit{of} your employment with the Company, to sign the Company’s standard Proprietary Information.\\
5 & In exchange for the severance benefits described in Section 2.3(a)(i) of my employment contract with the Company dated (…) (“Employment Contract”), which I acknowledge is sufficient \textit{consideration} to support this Agreement, & You also will be eligible for \textit{consideration} to receive long term incentive equity awards\\
\lspbottomrule
\end{tabularx}
\caption{Sample phrases with “term” and “consideration”.}
\label{tab:giampieri:5}
\end{table}

The examples reported in \tabref{tab:giampieri:5}  show that polysemy chaaracterize both corpora, although the EOL corpus produces fewer occurrences.

With regard to the polysemous word “term”, it is evident that line one of \tabref{tab:giampieri:5} contains instances where it means “duration” in both corpora, whereas it is a synonym of “word” in both the examples of line two. Finally, line three shows that “term” is a replacement for “condition”.

Also, the expression “in consideration” is polysemous in the EA corpus, as it signifies “as a counter-performance” (line four). In the EOL corpus, by contrast, this expression means “considering”; hence, it is not a controversial equivalent of “as a form of payment” (see line four). Conversely, in the fifth line, both terms are polysemous.

Complex prepositions hallmark the two corpora. If the expression “for the purpose* of” is queried in the EA corpus, 22 hits (0.57\%) are found, whereas only one occurrence is retrieved from the EOL corpus (0.03\%). The expression “pursuant to” is used 139 times in the EA corpus (3.6\%) against 33 (1.2\%) in the EOL corpus. The phrase “other than”, meaning “except”, shows 97 occurrences in the EA corpus (2.5\%), but only 23 (0.8\%) in the EOL corpus. Finally, the prepositional phrase “by virtue of” produces 4 hits (0.1\%) in the EA corpus, but no results in the EOL (see \tabref{tab:giampieri:6}  for phrases in context).

\begin{table}
\begin{tabularx}{\textwidth}{lQQ}
\lsptoprule
{No.} & {Sample phrase (EA)} & {Sample phrase (EOL)}\\
\midrule
1 & \textit{Pursuant to} the terms hereof, \textit{for the purposes of} entering judgment with respect to such award. & \textit{For the purpose of} calculating any benefits.\\
2 & Employment is terminated by the Company without Cause \textit{pursuant to} paragraph 16(a)(iv) below. & Severance payments made \textit{pursuant to} Section 9 or Section 10.\\
3 & Any attempted assignment, transfer, conveyance or other disposition (\textit{other than} as \textit{aforesaid}) of any interest. & You are not eligible for any severance benefits \textit{other than} those set forth in such agreement.\\
4 & The Employee will not copy or duplicate any of the \textit{aforementioned} documents. & The occurrence of one of the \textit{aforementioned} events.\\
5 & \textit{WHEREAS}, \textit{by virtue of} \textit{his} employment with Employer, Executive will become acquainted with certain confidential information. & {}-\\
\lspbottomrule
\end{tabularx}
\caption{Sample phrases with complex prepositions.}
\label{tab:giampieri:6}
\end{table}

As can be observed in line one and three of \tabref{tab:giampieri:6}, several complex prepositions appear in the same sentence of EAs (i.e., “pursuant to” and “for the purposes of” in line one, and “other than” and “aforesaid” in line three). Line four also features “aforementioned” in both corpora. In this respect, there are nine occurrences of “aforementioned” and “aforesaid” in the EA corpus (0.2\%), whereas only one instance of “aforementioned” in the EOL corpus (0.04\%) (line five). Finally, in line five, other \textit{legalese} traits come to the fore, such as sexist language (“his”), and the term “whereas”, which is an expression setting the contract premises or recitals.

In addition to the features explored above, there are other \textit{legalese} peculiarities in the two corpora which are not specifically quantifiable and, for this reason, are not reported in \tabref{tab:giampieri:1}. These elements are relevant for text comprehensibility and are, for example, repetitions and use of near synonyms; nominalization; long sentences, and syntactic discontinuities. Repetitions and near-synonyms, for example, characterize both employment agreements and employment offer letters. \tabref{tab:giampieri:7}  reports some sample phrases.

\begin{table}
\begin{tabularx}{\textwidth}{lQQ}
\lsptoprule
{No.} & {Sample phrase (EA)} & {Sample phrase (EOL)}\\
\midrule
1 & (A) Computer programs, systems and/or software; ideas, inventions, trademarks, business information, know-how, processes, improvements, designs, redesigns, discoveries and developments.\newline
(B) Trade or business secrets, including business opportunities, marketing or business diversification plans, business development and bidding techniques, methods and processes, financial data and the like (collectively, the “Protected Information”). & (A) Trade secrets, confidential and proprietary information.\newline
(B) Trade secrets, confidential business or technical information or know-how not generally known to the public.\\
2 & Employer, its Affiliates, and all of their agents, employees, officers and directors. & [\textit{Company name}] and its affiliates and subsidiaries.\\
3 & (A) Company policies, welfare or benefit plans, or other employment practices.\newline
(B) Plans, policies and arrangements. & The Company’s practices and policies.\\
\lspbottomrule
\end{tabularx}
\caption{Sample phrases with repetitions and near-synonyms.}
\label{tab:giampieri:7}
\end{table}

The clauses shown in line 1 of \tabref{tab:giampieri:7}  concern the prohibition for the employee to use trademarks or copyrighted materials during or after employment. To this aim, all the company’s proprietary information is listed. Line two focuses on the people working with or for the employer, while line three refers to the employer’s policies and customs. Although near-synonyms characterize both corpora, the lists mentioned in employment agreements is longer and more inclusive.

As discussed above, another interesting \textit{legalese} trait is nominalization. Both types of documents feature nominalization. For example, “make” is often followed by an object, such as “make an election” (instead of “elect” or, better, “choose”), “make payment” (instead of “pay”), “make changes” (instead of “change”), etc. These and other examples are reported in \tabref{tab:giampieri:8}.

\begin{table}
\begin{tabularx}{\textwidth}{lQQ}
\lsptoprule
{No.} & {Sample phrase (EA)} & {Sample phrase (EOL)}\\
\midrule
1 & (A) You shall advise the Company in writing whether or not you have \textit{made an election}.\newline
(B) The Corporation shall have no further obligation to \textit{make any payment}.\newline
(C) The Corporation shall not \textit{make any changes} in such plans or programs. & (A) We neither need nor desire to \textit{make any unauthorized use} of any intellectual property.\newline
(B) You will not need to \textit{make any cash payments} to cover such tax withholding.\newline
(C) Please \textit{make arrangements} with me to complete the necessary eligibility to work.\\
2 & Upon \textit{presentation} of such documentation thereof. & (A) Upon \textit{presentation} of appropriate documentation in accordance with the Company’s expense reimbursement policies.\newline
(B) You will need to present documentation.\\
3 & (A) Based on the \textit{achievement} of performance goals.\newline
(B) Based on achieving a combination of revenue, profit and other business performance targets. & (A) Based upon the \textit{achievement} of performance targets.\newline
(B) Subject to [\textit{Company name}] achieving certain stated performance goals.\\
\lspbottomrule
\end{tabularx}
\caption{Sample phrases with nominalization.}
\label{tab:giampieri:8}
\end{table}

In line two (B) of \tabref{tab:giampieri:8}, the plain expression “you will need to present documentation” can replace the nominalized phrase “upon presentation of such documentation”, which appears in both corpora. In the same way, the expression “based on the achievement of” could be simplified by writing “based on / subject to achieving” (line three, B).

Also long sentences (and syntactic discontinuity) characterize legal documents. In the current analysis, sentences are considered long when they contain more than 50 words (see \citealt{Williams2023}: 67–68). \citet{Williams2023}, in fact, defines as long the sentences which exceed 45 words \citep[67]{Williams2023}. The corpus of employment agreements produces more instances in this respect. One of them is reported in \tabref{tab:giampieri:9}.

\begin{table}
\begin{tabularx}{\textwidth}{lQQ}
\lsptoprule
{No.} & {Sample phrase (EA)} & {Sample phrase (EOL)}\\
\midrule
1 & Individuals who, \uline{as of the date \textit{hereof}}, constitute the Board (the “Incumbent Board”) cease for any reason to constitute at least a majority of the Board; provided, however, that any individual (\uline{\textit{other than} an individual whose initial assumption of office occurs as a result of an actual or threatened solicitation of proxies or consents by or on behalf of a Person other than the Board}) who becomes a director \textit{subsequent to} the date \textit{hereof} \uline{whose election or nomination for election by [\textit{Company name}]’s shareholders was approved by a vote of at least a majority of the directors then comprising the Incumbent Board} shall be considered as though such individual were a member of the Incumbent Board.
& 
In the event that:
(i) the Company consummates a change of control transaction, whereby fifty percent (50\%) or more of the voting stock of the Company changes ownership \textit{pursuant to} such transaction (a “Change of Control”); and
(ii) within twelve (12) months after the \textit{consummation} of a Change of Control, your employment with the Company is either (a) terminated by the Company or successor entity without Cause (\uline{as defined in Section 9  \textit{herein}}), or (b) terminated by you for Good Reason (\uline{as defined in and in accordance with the paragraph below}); and (iii) you sign, date, and deliver to the Company the Release and allow it to become effective; then you will receive the following as your \textit{sole} severance benefits (the “Change of Control Severance Benefits”): \\
\lspbottomrule
\end{tabularx}
\caption{Long sentences and syntactic discontinuities.}
\label{tab:giampieri:9}
\end{table}

\largerpage
Both the clauses of \tabref{tab:giampieri:9} are entitled “change of control”. The clause of the EA corpus is composed of 116 words; hence, it is extremely long and complex. Syntactic discontinuities are underlined and there are many \textit{legalese} features, such as “hereof”, “other than”, “subsequent to”, and “shall”. The clause sourced from the EOL corpus is even longer (126 words) and the one reported in \tabref{tab:giampieri:9} is only an extract of the entire clause. Other \textit{legalese} patterns are as follows: “pursuant to”, “herein”, nominalization (“the consummation of”), and “sole”. Although longer, the EOL clause contains fewer instances of syntactic discontinuity. Furthermore, the organization of the text is more straightforward. Sequencing and indexing, in fact, make the clause more comprehensible. Also, the tone is less formal than the one of the corresponding EA clause. The clause reported in \tabref{tab:giampieri:9}   is only an example of many similar (long) sentences of the clauses composing the EA and EOL corpora.

The overall tone of both employment agreements and employment offer letters is now scrutinized. \tabref{tab:giampieri:10}  lists a few sample sentences and phrases, where it is evident that the register used in employment letters is far more direct and, to some extent, plainer.

\begin{table}
\begin{tabularx}{\textwidth}{lQQ}
\lsptoprule
{No.} & {Sample phrase (EA)} & {Sample phrase (EOL)}\\
\midrule
1 & \textit{The Executive} \textit{shall be eligible to} receive a target annual cash bonus of forty-five percent (45\%) of his base salary under The New York Times  {Company 1991} Executive Cash Bonus Plan (the “Executive Cash Bonus Plan”). & (A) \textit{You will be eligible} for an annual bonus of up to 20\% of your base salary.\newline
(B) \textit{You may be eligible} to receive an annual discretionary bonus.\\
2 & \textit{In the event that Executive} \textit{accepts employment} with, or provides any services to (whether as a partner, consultant, joint venturer or otherwise), any person or entity which offers products or services that are \textit{competitive with any products or services offered by the Company} or with any products or services that Executive is aware the Company intends to offer,\newline
Executive shall be deemed to have voluntarily \textit{terminated employment} with the Company & \textit{You must not} engage in any work, paid or unpaid, that creates an actual conflict of interest with the Company.\\
\lspbottomrule
\end{tabularx}
\caption{Tone.}
\label{tab:giampieri:10}
\end{table}

The first line of \tabref{tab:giampieri:10} focuses on an annual bonus that the employee may receive. The tone used in employment letters is more informal and direct (e.g., “you will be eligible for an annual bonus” against “the executive shall be eligible to receive a target annual cash bonus”).

The second line sets the prohibition for the employee to engage in any activity which could compete with the employer’s. Also in this case, the sentence of the letter of employment is much shorter, more direct and to the point.

\section{Discussion}\label{sec:giampieri:5}

The analysis carried out in the sections above showed a recurrent and systematic use of \textit{legalese} in employment agreements. By contrast, employment offer letters are characterized by \textit{legalese} to a lesser extent. There are long sentences, complex prepositions, passive forms and nominalization, but \textit{legalese} features are less pervasive, texts are smoother and, hence, more understandable.

For example, despite the persistence of “shall” in many letters of employment, there is also a recurrent use of “will” and “must” (see \tabref{tab:giampieri:1} and 2). Passive forms and polysemous words characterize both corpora, but the EOL corpus produces fewer instances. Complex prepositions hallmark employment agreements and employment letters, but those in employment offer letters are fewer than the prepositions found in employment agreements (see \tabref{tab:giampieri:1}, where “other than” shows a frequency percentage value of 2.5\% in the EA corpus, but only 0.8\% in the EOL corpus). Phrases with near-synonyms and repetitive words are mentioned in both corpora, but in letters of employment they are far shorter (see \tabref{tab:giampieri:7}, for example). Sentences are long in the two corpora, but those in employment offer letters tend to be organized in bullet points and/or by sequencing. These writing strategies make EOLs more comprehensible and easier to read (see \tabref{tab:giampieri:9}). In addition, by considering the overall tone of employment offer letters, it is evident that they are more direct, if not almost colloquial, in comparison with employment agreements.

Such considerations can be taken into account when delivering intralingual translations and when addressing target texts which, to some extent, must be drawn up in plainer language.

\section{Conclusions and future research}\label{sec:giampieri:6}

This chapter was aimed at exploring the legal language used in employment agreements and employment offer letters. In particular, it wished to bring \textit{legalese} features to the surface, analyse them and investigate to what extent employment agreements and employment offer letters make use of \textit{legalese} patterning in native speaking contexts. In this way, it also investigated possible corpus-assisted applications in (intralingual) translation studies.

This chapter contributed to highlighting the fact that employment agreements are drafted in a formal and archaic style, despite being addressed to both native and non-native speakers of English. They feature \textit{legalese} patterns, which are used systematically throughout the documents. By contrast, in employment offer letters, alternative (and plainer) linguistic solutions are proposed. Long sentences, syntactic discontinuities, nominalization, or complex prepositions are still used by the drafters of employment contracts and letters of employment, but the pedantic aspect of the language of the law seems to be less pervasive. As a matter of fact, plainer linguistic solutions are proposed in employment offer letters, such as shorter all-comprehensive lists of near-synonyms; a reduced recurrence of “shall” in favour of a more systematic use of “will”; fewer deixis patterns and simpler passive forms. Also, despite the challenging sentence length, letters of employment are organized in bullet points or in sequences, thus allowing readers to follow the text content easily.

Differently from what could be expected, the findings of this chapter highlight that employment agreements and employment offer letters do not differ in terms of the duties or positions offered to the recipient, as they are both addressed to a variety of employees, including executives. Hence, the reason(s) why companies decide to resort to one type of document instead of the other might be further explored by researchers.

There are many implications in the field of translation studies. Users can explore either comparable or parallel corpora (\citealt{FantinuoliZanettin2015}; \citealt{Resi2024}) in order to understand and verify language nuances and make informed decisions as to the best translation candidates in a given context. Translators could also analyse collocations in different languages and retrieve topical vocabulary, as well as multi-words and lexical phrases. For example, as the language of the law is particularly characterized by formulae and fixed language expressions \citep{Giampieri2025}, translators could explore such patterns thanks to corpus consultation. Most of the corpus platforms and corpus software solutions available today contain features that allow users to carry out language investigations (\citealt{chapters/5-Cuchet}; \citealt{Resi2024}). At their disposal, translators have collocation listing options, as well as key-word-in-context functions and lemmatization.

By creating a similar corpus composed of Italian \textit{contratti di impiego} (back-translation: “employment contract”), it is possible to notice words and expressions which translate most of the terms analyzed above. Examples of such phrases are \textit{premesso che} (back-translation: “whereas”); \textit{orario di lavoro} (back-translation: “working time”); \textit{traslocare} or \textit{trasferirsi} (back-translations: “move to” or “transfer to”); \textit{recedere dal contratto di impiego} (back-translation: “terminate employment”); \textit{avere diritto a} (back-translation: “be entitled to”); \textit{rimborsare spese} (back-translation “to reimburse expenses”), just to mention a few. An Italian corpus could be composed by searching for templates or drafts of \textit{contratto di impiego} in Italian domains (search query: \textit{“modello contratto di impiego” site:.it}). In this way, sample contracts are downloaded and their language and terminology can be analyzed by using the AntConc software solution.

The limit of this study lies in the corpora taken into consideration. An analysis of the contracts drawn up in English as a \textit{lingua franca}, for example, could have yielded more comprehensive, or perhaps different, results. In this respect, it could be interesting to explore and verify to what extent employment offer letters and/or employment agreements are used in non-native speaking contexts and which type of language (i.e., plainer or more intricate) and register or tone are applied. This is something that future research might wish to explore. Also, analysing the drafting style of other legal documents written in English as a first or second language could be another avenue to explore.

The implications of this study are manifold, as this chapter provides evidence that corpus consultation helps deliver intralingual translations \citep{Jakobson1959}. In other words, corpora can be used in the translation process to deliver target texts drafted in a plainer legal style. Hence, corpus consultation could be fostered for the benefit of both students in Translation Studies and professional translators. Further analyses could explore these aspects and propose the application of \textit{ad hoc} corpora in (intralingual) translator training and/or (intralingual) translation practice. The principles and methodologies followed in this chapter can be replicated in other fields of study, thus fostering interdisciplinarity in empirical translation research.

\sloppy\printbibliography[heading=subbibliography,notkeyword=this]
\end{document}
